% !TeX root = main.tex
\lecture{7}{Mon 10 Feb 2026 12:00}{Electrostatics VII: Electric Dipoles}

\section{Using Potential to Determine E}
Consider two charges $Q$ with distance $L$ between them. The potential at some generic $x$ is:

\[
    V(x) = \frac{Q}{4 \pi \epsilon_0 x} + \frac{Q}{4 \pi \epsilon_0 (L  -x)}
\]
\[
    \pmb{E} = - \nabla V = - \frac{dV}{dx} \hat{\pmb{x}}
\]
\[
    = \left(\frac{Q}{4 \pi \epsilon_0 x^2} - \frac{Q}{4 \pi \epsilon_0 (L - x)^2}\right) \hat{\pmb{x}}
\]

At the middle point $x = L/2$:
\[
    V(x = L/2) = \frac{Q}{4 \pi \epsilon_0 L/2} + \frac{Q}{4 \pi \epsilon_0 (L - L/2)}
\]
\[
    \pmb{E}(x = L/2) = \frac{Q}{4 \pi \epsilon_0 L^2/4} - \frac{Q}{4 \pi \epsilon_0 (L^2 / 4)} = 0
\]

\section{Potential from a Ring}
\begin{figure}[H]
    \centering
    \includegraphics[width=0.75\textwidth]{figures/lec07-01.png}
     \caption{}
\end{figure}

A ring-shaped conductor of radius $a$ carries a total charge $Q$ uniformly distributed around it. Find the potential, and hence the electric field, at point $P$.

Consider some small element of charge $\delta Q$, at distance $r = \sqrt{a^2 + x^2}$ from $P$.

\[
    \delta V = \frac{\delta Q}{4 \pi \epsilon_0 r} = \frac{\delta Q}{4 \pi \epsilon_0 (a^2 + x^2)^{1/2}}
\]

We can then integrate, noting that $a, x$ are both constants:
\[
    V = \frac{1}{4 \pi \epsilon_0 (a^2 + x^2)^{1/2}} \oint dQ
\]
\[
    = \frac{Q}{4 \pi \epsilon_0 (a^2 + x^2)^{1/2}}
\]

And then to find $\pmb{E}$, noting that $\pmb{E} = - \nabla V$ (and that the derivatives in the z, y directions are zero):
\[
    \pmb{E} = - \frac{\partial V}{\partial x} \pmb{\hat{x}} = - \frac{Q}{4 \pi \epsilon_0} \left(- \frac{1}{2} 2x (a^2 + x^2)^{-3/2}\right)
\]
\[
    = \frac{Qx}{4 \pi \epsilon_0 \left(a^2 + x^2\right)^{-3/2}} \hat{\pmb{x}}
\]

\section{Electric Dipoles}
Consider two charges $+q$ and $-q$, separated by a distance $a$. This forms an electric dipole. This has an associated quantity given the electric dipole moment:
\[
    \pmb{p} = q \pmb{a}
\]
We define this vector to point against the \pmb{E}-field, i.e. pointing from the negative towards the positive charge.

Many molecules naturally form dipoles. Because water is not vertically symmetric, it causes a difference between the effective positive charge from the two hydrogens and the net negative charge from the oxygen. 

This makes it an excellent solvent and is used in heating food via microwaves. Because the distance is very small, it gives a very small value of $p = 6.1 \times 10^{-30}$Cm per molecule.

\begin{figure}[H]
    \centering
    \includegraphics[width=0.5\textwidth]{figures/lec07-02.png}
     \caption{}
\end{figure}

For a dipole, the potential is given by:
\[
    V_p \approx \frac{p \cos \theta}{4 \pi \epsilon_0 r^2}
\]
\[
    \pmb{E} \approx \frac{2 p \cos \theta}{4 \pi \epsilon_0 r^3} \hat{\pmb{r}} + \frac{p \sin \theta}{4 \pi \epsilon_0 r^3} \hat{\pmb{\theta}}
\]
For $r \gg a$ and where $\theta$ is defined as the angle between $\pmb{r}$ and the positive charge.s






